\uuid{8yUr}
\exo7id{7453}
\titre{exo7 7453}
\auteur{mourougane}
\organisation{exo7}
\datecreate{2021-08-10}
\isIndication{false}
\isCorrection{false}
\chapitre{Géométrie affine dans le plan et dans l'espace}
\sousChapitre{Géométrie affine dans le plan et dans l'espace}
\cpge{mpsi-mp2i.espaces-vectoriels.d, pcsi.espaces-vectoriels.d, ptsi.espaces-vectoriels.d, mpsi-mp2i.sous-espaces-affines, ptsi.espaces-vectoriels.b, bcpst1.geometrie.a, bcpst1.geometrie.d, mp-mpi.endomorphismes-euclidien.c, pc.endomorphismes-euclidiens.a, psi.endomorphismes-euclidien.a, pt.isometries-euclidien.a, mp-mpi.endomorphismes-euclidien.d, pc.endomorphismes-euclidiens.c, psi.endomorphismes-euclidien.d, pt.isometries-euclidien.c, mp-mpi.endomorphismes-euclidien.e, psi.endomorphismes-euclidien.e, pt.isometries-euclidien.d, bcpst2.espaces-vectoriels.a}
\module{Géométrie}
\niveau{L2}
\difficulte{}

\contenu{
\texte{
Dans un espace vectoriel euclidien, deux sous-espaces vectoriels non nuls orthogonaux sont-ils toujours en somme directe ?
}
\begin{enumerate}
    \item \question{Montrer en utilisant la forme réduite que le déterminant d'une isométrie
est $(-1)^{codim E_1}$ où $E_1$ est l'espace propre de valeur propre $1$.}
    \item \question{Rappeler toutes les isométries du plan euclidien et de l'espace
euclidien de dimension 3, en précisant leur
partie linéaire, leur point fixe, leur axe, leur composante à point fixe et leur composante de glissement.}
    \item \question{Rappeler la construction de l'axe radical de deux cercles non sécants.}
    \item \question{Déterminer le groupe d'isométries d'un segment dans le plan euclidien.}
\end{enumerate}
}
